\section{Introduction}
\label{sec:introduction}

\paragraph{Scientific work is collaborative.}
Language-model agents now perform substantial parts of the scientific
workflow. They search literature, propose hypotheses, write and execute
analysis code, and assemble results into reports
\citep{lu2024aiscientist,yamada2025aiscientistv2,gottweis2025coscientist,mitchener2025kosmos}.
The benchmarks used to measure this progress have followed the same
trajectory, moving from recall and reasoning questions toward executable
analysis tasks drawn from real studies
\citep{majumder2024discoverybench,mitchener2025bixbench,kon2025expbench,ye2025replicationbench}.
Yet in practice these systems function as co-scientists rather than
autonomous investigators, and the case has been made that they are not
architecturally built for unsupervised discovery \citep{bisht2026notbuilt}.
A deployed scientific agent sits inside a collaboration: someone designed
the experiment, ran the instrument, and holds the context that determines
whether an analysis is valid. That context is frequently undocumented. Which
plates were processed on a different day, why three samples were dropped,
whether a covariate is confounded with batch, which normalization the lab
actually applies---these facts live with the researcher, not in the data
directory. The agent's scientific output therefore depends not only on its
own reasoning but on an exchange with a particular human who may explain
readily, answer only what was literally asked, or misremember a detail.

\paragraph{The evaluation gap.}
Most scientific agent benchmarks specify each task completely and
uniformly, which measures capability under ideal information but says
nothing about capability under realistic information transfer. A recent line
of interactive benchmarks closes part of this gap by requiring agents to
recognize underspecification and ask: HiL-Bench measures whether coding
agents escalate at the right moments \citep{trinh2026hilbench},
SciConvBench evaluates multi-turn clarification for ill-posed computational
science requests \citep{somasekharan2026sciconvbench}, and work in
tool-agent settings evaluates interaction with a simulated user under
domain rules \citep{yao2024taubench,chen2026appworldul,ta2026regretbench}.
These benchmarks establish that asking is hard. They do not, however, vary
\emph{whom} the agent is asking. The simulated collaborator is typically a
single cooperative interlocutor with one disclosure policy, so the
collaborator is a fixed property of the environment rather than an
experimental condition. The closest departures from this pattern vary
designed behavioral axes---non-collaborative user simulators degrade tool
agents through impatience, digression, and incomplete utterances
\citep{shim2025noncollaborative}, and evolved persona policies induce
realistic behavioral variation in otherwise cooperative simulators
\citep{chopra2026ppol}---but do so in retail, airline, and customer
service domains where there is no scientific ground truth to get wrong.
The result is that three distinct sources of difficulty remain entangled in
current evaluation: the intrinsic difficulty of the scientific task, the
capability of the agent, and the character of the collaborator the agent
happens to be working with. Reporting a single score against a single
default collaborator silently averages over the third.

\paragraph{Our approach.}
This paper introduces \benchmark, a benchmark that holds the scientific
task fixed and varies the simulated researcher. We convert 200 tasks from
five scientific agent benchmarks---DiscoveryBench, BixBench, EXP-Bench,
SciConvBench, and ASI-Bench---into collaboration tasks, preserving each
source's native deliverable and grading rules rather than reducing them to
a common format. An actor agent must produce that native deliverable while
interacting with a simulated researcher through a shared review protocol.
The researcher is conditioned on a complete rendered profile drawn from
MatrAIx, a population-scale simulated-user infrastructure whose persona
records span 1,290 categorical dimensions covering background, psychology,
capability, and behavior \citep{li2026matraix}. Our renderer emits the
populated attributes as an explicit attribute--value listing grouped by
category, so a condition is a fully auditable set of stated attributes rather
than authored prose; a single profile typically contributes on the order of two
hundred such attributes. The profile is injected verbatim into the researcher
simulator's system prompt only, without truncation or projection onto the
task domain, so the actor observes personality solely through conversation.
Every task is also run with the identical protocol and the
identity block removed, giving a paired persona-off control that separates
collaborator variation from task difficulty. Evaluator references stay
private throughout, and each task--condition cell is repeated across
independent episodes so that persona effects can be distinguished from
ordinary run-to-run variation. Because the same task is scored under
multiple collaborators with a shared baseline, differences are attributable
to the collaborator rather than to task selection.

Upstream validation of this persona corpus establishes that conditioning
reliably induces declared behavior---91.5\% of trials in a controlled
adherence study---but does so in Survey, Chatbot, Web, and App environments
over a task library concentrated in commerce, software, finance, and
healthcare \citep{li2026matraix}. Scientific research collaboration is not
among those environments, and product feedback rather than scientific
correctness is the measured outcome. \benchmark supplies that missing
environment and asks the question it enables: whether persona-induced
behavioral variation propagates into graded scientific error.

\paragraph{What we measure.}
Evaluation is deliberately two-sided. Outcome metrics use each source
benchmark's native grader, so a persona effect registers as a change in
scientific correctness and not merely as a change in conversational tone.
Interaction metrics---review rounds, revisions, questions asked,
clarification coverage, message length, and language adaptation---are
recorded separately, because a collaborator can change an agent's behavior
substantially without changing its score, and can change its score without
changing observable behavior. Keeping these separate is what allows a
failure under a difficult collaborator to be diagnosed as a failure to
obtain available information rather than assumed to be one.
\todo{Insert the principal quantitative findings once the repeated-run
analysis over the full persona panel is complete: direction and magnitude of
persona effects with confidence intervals, which benchmarks and fields are
most sensitive, and whether behavioral divergence tracks outcome divergence.
Pilot runs showed score changes in both directions of substantial magnitude
relative to the shared baseline, but that panel was outcome-selected and is
reported as a selection stage, not as an effect estimate.}

\paragraph{Contributions.}
\begin{itemize}[leftmargin=1.5em]
  \item \textbf{A persona-conditioned scientific collaboration benchmark.}
  We convert 200 tasks from five scientific agent benchmarks into
  interactive collaboration tasks under a source-preserving contract that
  retains native deliverables, graders, and private references
  (\secref{sec:benchmark}).
  \item \textbf{An experimental design that isolates the collaborator.}
  Full-profile researcher conditioning, a paired persona-off control on the
  identical protocol, and repeated independent episodes per cell yield
  paired comparisons in which the researcher is the manipulated variable
  (\secref{sec:experimental-setup}). We document the two-stage profile
  selection, including its outcome-dependent final stage, rather than
  presenting the panel as population-representative.
  \item \textbf{An empirical characterization of collaborator sensitivity.}
  We report outcome and interaction effects with repeat-based uncertainty,
  their dependence on task format and scientific field, and
  evidence-linked conversation cases tracing profile wording to researcher
  behavior, actor decisions, and graded outcomes
  (\secref{sec:results},~\secref{sec:discussion}). Completed experiments
  are reported separately from proposed extensions.
\end{itemize}

\noindent
\secref{sec:related-work} positions this work against scientific agent
benchmarks, interactive and simulated-user evaluation, and persona
simulation. \secref{sec:benchmark} defines the benchmark,
\secref{sec:experimental-setup} the experimental protocol,
\secref{sec:results} the results, and \secref{sec:discussion} mechanism
evidence from conversation traces. \secref{sec:limitations} states what
simulated collaborators and an outcome-selected panel cannot establish.
